\documentclass[11pt,a4paper]{article}
\usepackage[utf8]{inputenc}
\usepackage[T1]{fontenc}
\usepackage[brazilian]{babel}
\usepackage[margin=2.3cm]{geometry}
\usepackage{booktabs,array,tabularx}
\usepackage{xcolor}
\usepackage{pgfplots}
\pgfplotsset{compat=1.17}
\PassOptionsToPackage{hyphens}{url}
\usepackage[hidelinks]{hyperref}
\usepackage{enumitem}
\setlist{nosep,leftmargin=1.4em}

\definecolor{luna}{HTML}{3B6EA8}
\definecolor{sol}{HTML}{D08A1E}
\definecolor{cache}{HTML}{B9C6D6}

\title{Sweatshop no Codex: primeira leva\\\large gpt-6-luna (max) no stage 2, gpt-6-sol (high) no stage 3}
\author{Foreman, sessão \texttt{sweatshop/2026-09-24-1506}}
\date{24 de setembro de 2026}

\begin{document}
\maketitle

\section*{Resumo}

Primeira execução do driver \texttt{sweatshop-codex.ps1} no SynchroNice, com a
lineup nova: \textbf{gpt-6-luna max} implementando e \textbf{gpt-6-sol high}
revisando. O run foi encerrado a pedido após o primeiro ticket mergeado.

\begin{itemize}
  \item \textbf{Produzido:} 1 ticket (SYN-010, PC: Pasta padrão e engrenagem),
        mergeado na sessão após \textbf{3 ciclos} de implementação e revisão.
        PR de sessão: \href{https://github.com/Laginho/SynchroNice/pull/4}{\#4}.
  \item \textbf{Custo do ticket:} 95 min de relógio, 20,8\,M tokens de entrada
        (96,8\,\% em cache) e 160\,k de saída (72\,\% raciocínio).
  \item \textbf{Integração Codex:} um bug bloqueante no \texttt{ticket-flow}
        (guarda de modelo), corrigido durante o run; a integração em si
        (commit, push, gate, merge na sessão, PR) funcionou.
  \item \textbf{Veredito provisório:} o luna produz código correto e respeita o
        TDD, mas é mais lento que a lineup Claude. O gargalo maior foi o
        \textbf{sol}: revisões incrementais e com aumento de escopo custaram dois
        ciclos extras (52\,\% do tempo do ticket). Com $n=1$, nada disso é
        conclusivo.
\end{itemize}

\section{Linha do tempo}

\begin{tabularx}{\linewidth}{@{}lX@{}}
\toprule
15:06 & Primeiro lançamento. Quatro estágios (SYN-010 e SYN-012, duas tentativas cada)
        recusam trabalhar em cerca de 1 min cada: o Codex não informa à sessão o
        nome do modelo, e a guarda \emph{``Wrong model, no work''} do
        \texttt{ticket-flow} manda parar. Os dois tickets ficam \texttt{blocked}.\\
15:11 & Driver para (``nothing runnable''). Correção no skill
        (Agent-Skills \texttt{a8cdadb}): no Codex, confiar no \path{--model} do
        lançamento. Tickets de volta a \texttt{to-implement}.\\
15:14 & Segundo lançamento. SYN-010, implementação 1.\\
15:54 & \texttt{to-review} após 40 min. Revisão 1 (6 min): \textbf{reaberto}.\\
16:15 & Implementação 2 (15 min). Revisão 2 (3 min): \textbf{reaberto}.\\
16:43 & Implementação 3 (24 min). Revisão 3 (7 min): \textbf{mergeado}.\\
16:50 & Encerramento pedido: o estágio seguinte (SYN-011) é interrompido nos
        primeiros segundos e o driver fecha pelo próprio \texttt{finally}
        (árvore limpa, PR atualizado, volta a \texttt{main}).\\
\bottomrule
\end{tabularx}

\section{Métricas por estágio (SYN-010)}

Tokens do evento \texttt{turn.completed} do \path{codex exec --json}.
``Cmds'' é o número de comandos de shell que a sessão executou.

\medskip
\noindent{\small
\begin{tabularx}{\linewidth}{@{}l l r r r r r X@{}}
\toprule
Estágio & Modelo & Min & Cmds & Entrada & Cache & Saída & Desfecho\\
\midrule
Impl.\ 1   & luna & 40 & 192 & 10,61\,M & 97,9\,\% & 75,4\,k & to-review\\
Revisão 1  & sol  &  6 &  48 &  1,04\,M & 94,2\,\% &  5,5\,k & reaberto (crit.\ 5)\\
Impl.\ 2   & luna & 15 &  84 &  2,50\,M & 95,7\,\% & 28,5\,k & to-review\\
Revisão 2  & sol  &  3 &  32 &  0,74\,M & 93,5\,\% &  5,5\,k & reaberto (crit.\ 4, 5)\\
Impl.\ 3   & luna & 24 &  86 &  3,60\,M & 96,4\,\% & 37,8\,k & to-review\\
Revisão 3  & sol  &  7 & 118 &  2,28\,M & 96,0\,\% &  7,0\,k & mergeado\\
\midrule
\multicolumn{2}{@{}l}{Total} & 95 & 560 & 20,76\,M & 96,8\,\% & 159,6\,k & \\
\multicolumn{2}{@{}l}{Recusas (bug)} & 4 & 28 & 0,42\,M & 77,9\,\% & 10,5\,k & 4 estágios perdidos\\
\bottomrule
\end{tabularx}}

\bigskip
\begin{center}
\begin{tikzpicture}
\begin{axis}[
  ybar stacked, bar width=18pt, width=0.92\linewidth, height=6cm,
  ylabel={Tokens de entrada (M)}, ymin=0,
  symbolic x coords={Impl 1,Rev 1,Impl 2,Rev 2,Impl 3,Rev 3},
  xtick=data, legend style={at={(0.98,0.95)},anchor=north east,draw=none},
  nodes near coords align={vertical},
  axis lines*=left, ymajorgrids, grid style={gray!25}]
\addplot[fill=cache,draw=none] coordinates
  {(Impl 1,10.386) (Rev 1,0.977) (Impl 2,2.390) (Rev 2,0.695) (Impl 3,3.471) (Rev 3,2.185)};
\addplot[fill=luna,draw=none] coordinates
  {(Impl 1,0.222) (Rev 1,0.060) (Impl 2,0.109) (Rev 2,0.049) (Impl 3,0.129) (Rev 3,0.091)};
\legend{em cache, sem cache}
\end{axis}
\end{tikzpicture}
\end{center}

\paragraph{Onde foi o custo.} Os ciclos 2 e 3, que só existiram por causa das
reaberturas, somaram 49 min (52\,\% do relógio), 9,1\,M de entrada (44\,\%) e
79\,k de saída (49\,\%). A primeira implementação sozinha foi metade da entrada
do ticket; parte disso foi a tentativa de verificação visual pelo Computer Use
(ver seção~\ref{sec:integracao}).

\section{Comparação com a lineup Claude}

Referência: os 13 tickets mergeados nos runs anteriores com Sonnet 5 no stage 2
e Opus 5.5 no stage 3 (\texttt{run-log.md}). Não há tokens para eles: o driver
Claude grava a saída em texto, sem contagem.

\medskip
\noindent{\small
\begin{tabularx}{\linewidth}{@{}X r r@{}}
\toprule
& Claude (13 tickets) & Codex (SYN-010)\\
\midrule
Minutos por ticket mergeado, mediana (faixa) & 42 (7--89) & 95\\
Tickets reabertos & 5 de 13 & 1 de 1\\
Reaberturas no mesmo ticket, máximo & 1 & 2\\
Minutos por implementação que chegou a \texttt{to-review}, média & 14,7 & 26,3\\
Minutos por revisão, média & 9,2 & 5,3\\
\bottomrule
\end{tabularx}}

\medskip
SYN-010 foi mais lento que qualquer ticket da lineup Claude. Ressalvas: é um
ticket largo (backend Rust, UI e casos de borda do Windows, 12 arquivos e
$+406/-73$ na sessão), a comparação é com uma amostra de outro tamanho, e parte
do tempo da implementação 1 foi gasta com o ambiente, não com o modelo.

\section{Observações sobre os modelos}

\subsection*{gpt-6-luna (max), implementação}
\begin{itemize}
  \item \textbf{TDD respeitado:} em cada ciclo, commit só de teste, vermelho,
        antes do código; a revisão 3 aprovou os dois critérios reabertos com
        testes de regressão nos seams pedidos (37 $\rightarrow$ 39 testes Rust).
  \item \textbf{Uma infração pequena:} \texttt{0a88f83 test: format default
        folder assertion} mexe num teste depois do commit \texttt{feat}. É só
        formatação, mas o \texttt{ticket-flow} proíbe, e nenhuma das três
        revisões apontou.
  \item \textbf{Bom senso operacional:} ao verificar a tela, usou um perfil
        temporário, desfez a entrada de autostart que a própria instância criou
        e registrou o que não conseguiu verificar como débito humano, em vez de
        declarar sucesso.
  \item \textbf{Caro:} 192 comandos e 10,6\,M de entrada na primeira passada.
\end{itemize}

\subsection*{gpt-6-sol (high), revisão}
\begin{itemize}
  \item \textbf{Revisão incremental.} Ao que tudo indica, os três defeitos apontados já estavam no
        código na primeira passada, mas foram listados um por revisão (1, depois
        2). Cada achado adiado custou um ciclo de cerca de 20 min e 3--4\,M
        tokens. O Opus nunca reabriu o mesmo ticket duas vezes.
  \item \textbf{Aumento de escopo.} Nenhum dos três achados estava nos
        critérios originais. A revisão reescreveu o texto dos critérios 4 e 5
        (engrenagem acessível na tela de QR; Pasta gravada funcionando sem
        Downloads resolvível; arquivos antigos na Pasta original). O primeiro é
        uma regressão real introduzida pela feature; os outros dois são
        requisitos novos, que pelo fluxo iriam para o stage 1 ou para um
        \texttt{CLEAN-*}.
  \item \textbf{Deixou passar} a infração de TDD acima.
  \item \textbf{Barato por passada:} 3--7 min, 0,7--2,3\,M de entrada.
\end{itemize}

\section{Integração com o Codex}\label{sec:integracao}
\begin{itemize}
  \item \textbf{Bug corrigido:} a guarda de modelo do \texttt{ticket-flow}
        pressupunha que a sessão conhece o próprio modelo. No Codex isso não é
        verdade, e todo estágio parava. Correção em Agent-Skills
        \texttt{a8cdadb}.
  \item \textbf{Funcionou:} commit e push de dentro do estágio, gate completo
        (fmt, clippy, testes Rust, \texttt{tsc}, Gradle), merge na branch de
        sessão, publicação do PR, contagem de tentativas, tratamento de erro de
        API.
  \item \textbf{Computer Use não funciona nesta máquina:} a captura de janela
        falha em \texttt{SetIs\allowbreak Border\allowbreak Required} (\emph{No such interface
        supported}). Toda verificação visual virou débito humano.
  \item \textbf{Sandbox:} os estágios rodam com
        \path{--dangerously-bypass-approvals-and-sandbox}. A regra de revisão
        automática do Codex barrou, corretamente, a exclusão de
        \path{%APPDATA%\dev.synchronice.pc}.
  \item \textbf{Sem parada limpa:} o driver não tem sinal de ``pare após o
        estágio atual''; o encerramento pedido foi feito interrompendo o
        estágio seguinte duas vezes, o que o driver lê como queda de API.
\end{itemize}

\section{Recomendações}
\begin{enumerate}
  \item \textbf{Subir o sol para max na próxima leva.} A revisão custa cerca de
        um terço de um ciclo de retrabalho; se o esforço maior evitar uma
        reabertura a cada dois tickets, já se paga.
  \item \textbf{Instruir contra aumento de escopo} no stage 3 do
        \texttt{ticket-flow}: revisar contra os critérios escritos, listar
        \emph{todos} os achados numa passada e mandar requisito novo para
        \texttt{CLEAN-*}. Sem isso, mais esforço pode gerar mais reaberturas.
  \item \textbf{Manter o luna max} até haver amostra: reavaliar depois de
        SYN-011 e SYN-012.
  \item \textbf{Adicionar ao driver um arquivo de parada} (terminar após o
        estágio atual), para não depender de matar processos.
\end{enumerate}

\section{Débito humano}
\begin{itemize}
  \item Inspeção visual do chat, painel, seletor nativo, tela inicial sem
        diálogo e menu da bandeja (Computer Use falhou).
  \item Confirmar a engrenagem e o painel sobre a tela de QR e no chat após a
        mudança do cabeçalho.
  \item Pastas deixadas pela inspeção, para limpar à mão (confira antes: as
        duas primeiras podem ser as do app real, com o pareamento do celular):
        \begin{itemize}
          \item \path{%APPDATA%\dev.synchronice.pc}
          \item \path{%LOCALAPPDATA%\dev.synchronice.pc}
          \item \path{%TEMP%\SynchroNice-SYN-010}
        \end{itemize}
\end{itemize}

\end{document}
